% ============================================================================
% MA 551 Computational Statistics -- Final Project
% Empirical Analysis of Saddle-Point-Escaping Optimizers (SPGD)
%
% Self-contained arXiv-preprint-style report.  Compiles with pdflatex on
% Overleaf with no extra .sty files; just upload report.tex, references.bib,
% and the figures/ directory.
%
% Build:
%   pdflatex report
%   bibtex   report
%   pdflatex report
%   pdflatex report
% ============================================================================

\documentclass[11pt]{article}

% ---- Page geometry / arXiv-preprint look -----------------------------------
\usepackage[margin=1in,headheight=14pt]{geometry}
\usepackage{microtype}
\usepackage{lmodern}
\usepackage[T1]{fontenc}
\usepackage[utf8]{inputenc}

% ---- Math, graphics, tables ------------------------------------------------
\usepackage{amsmath,amssymb,amsthm}
\usepackage{graphicx}
\usepackage{booktabs}
\usepackage{multirow}
\usepackage{array}
\usepackage{enumitem}
\usepackage[font=small,labelfont=bf]{caption}
\usepackage{subcaption}
\usepackage{xcolor}
% ---- Float placement -------------------------------------------------------
% `float`    -> the [H] specifier ("HERE, no really")
% `placeins` -> \FloatBarrier prevents floats from crossing subsection ends
\usepackage{float}
\usepackage{placeins}  % no [section] option: don't auto-insert \FloatBarrier
                       % at every section (that was causing big trailing
                       % whitespace when a section is shorter than the
                       % deferred figure).
% Be lenient with float-page packing so figures don't get postponed
% and leave whitespace behind.
\renewcommand{\topfraction}{0.95}
\renewcommand{\bottomfraction}{0.95}
\renewcommand{\textfraction}{0.05}
\renewcommand{\floatpagefraction}{0.6}
\setcounter{topnumber}{4}
\setcounter{bottomnumber}{3}
\setcounter{totalnumber}{6}

% Vertically justify pages: any slack is absorbed into the existing
% rubber spacing around floats / paragraphs (\textfloatsep, \intextsep,
% \parskip), which spreads gaps invisibly instead of letting a single
% deferred float leave a large block of trailing whitespace.
\flushbottom

% Tight nominal spacing around floats (default \article spacing is
% generous), with enough rubber stretch (plus...) that \flushbottom can
% absorb page slack across multiple gaps instead of dumping it at the
% bottom of the page.
\setlength{\textfloatsep}{8pt plus 6pt minus 2pt}
\setlength{\intextsep}{8pt plus 6pt minus 2pt}
\setlength{\floatsep}{8pt plus 6pt minus 2pt}
\setlength{\abovecaptionskip}{4pt}
\setlength{\belowcaptionskip}{2pt}

% ---- Algorithms ------------------------------------------------------------
\usepackage{algorithm}
\usepackage{algpseudocode}

% ---- Bibliography (author-year, arXiv-style) ------------------------------
\usepackage[numbers,sort&compress]{natbib}

% ---- Hyperlinks (load last) ------------------------------------------------
\usepackage[colorlinks=true,
            linkcolor=blue!50!black,
            citecolor=blue!50!black,
            urlcolor=blue!50!black]{hyperref}

% ---- Fancy header (light) --------------------------------------------------
\usepackage{fancyhdr}
\pagestyle{fancy}
\fancyhf{}
\renewcommand{\headrulewidth}{0pt}
\fancyhead[L]{\small\textsc{SPGD on Non-Convex ML Loss Landscapes}}
\fancyhead[R]{\small \thepage}

% ---- Title block (arXiv-preprint style) -----------------------------------
\title{\bfseries \LARGE Does Steepest Selection Help?\\
       An Empirical Study of Saddle-Point-Escaping Optimizers\\
       on Non-Convex Machine Learning Loss Landscapes}

\author{%
  \textbf{Shyam Patadia}\\
  Worcester Polytechnic Institute (WPI)\\
  \texttt{spatadia@wpi.edu}\\[2pt]
  \small MA 551: Computational Statistics\\[2pt]
  \small Code: \url{https://github.com/shyampatadia/spgd_research}
}

\date{\today}

% ---- Convenience macros ----------------------------------------------------
\newcommand{\np}{\ensuremath{N_P}}
\newcommand{\iterp}{\ensuremath{\textit{IterP}}}
\newcommand{\amp}{\ensuremath{\textit{Amp}}}
\newcommand{\R}{\mathbb{R}}
\newcommand{\E}{\mathbb{E}}

\begin{document}
\maketitle
\thispagestyle{empty}

% ============================================================================
\begin{abstract}
\noindent
Training a modern machine learning model is, mathematically, the
minimization of a high-dimensional non-convex loss function. In such
landscapes, saddle points dominate the population of critical points and
constitute the chief obstacle to first-order optimization. \emph{Steepest
Perturbed Gradient Descent} (SPGD) \citep{vahedi2024spgd} is a
candidate-selection rule on top of gradient descent, motivated by
saddle escape: every \iterp{} iterations, sample \np{} candidate
solutions in a uniform neighborhood of the current iterate and accept
the candidate that yields the steepest loss decrease (the paper's
$\leq$ acceptance rule). SPGD has only been
evaluated on smooth mathematical benchmarks and a small engineering
problem; its behavior on neural-network loss landscapes has not been
characterized. This work provides the first systematic empirical study
of SPGD on machine-learning workloads. We implement SPGD as a PyTorch
\verb|Optimizer|, compare it against SGD, Adam, perturbed gradient
descent (PGD) \citep{jin2017escape}, and a same-mechanism
random-perturbation control (RPGD) on four discriminating testbeds
in the main study --- non-convex synthetic benchmarks, OpenML-CC18
heterogeneous tabular data, a CIFAR-10/ResNet-18 anchor experiment,
and a MovieLens-100K non-convex matrix-completion benchmark --- with
two further diagnostic experiments (Two Moons visualisation and an
$\np \times \iterp$ hyperparameter ablation) reported in the
appendix. We measure
saddle-point behavior directly via stagnation episodes
($\|\nabla f\| < \varepsilon$). Our key empirical finding is that
SPGD's steepest-selection rule provides a measurable advantage over
random perturbation across both synthetic and real-data regimes (mean
$+3.78$ percentage-point test-accuracy gap on CIFAR-10/ResNet-18,
paired across all seeds), but the gap relative to vanilla SGD is small
when no overt plateau is encountered, and Adam dominates absolute
accuracy on well-conditioned landscapes. We conclude that the value of
SPGD lies in its \emph{selection rule}, not in perturbation alone, and
sketch when this advantage is operationally relevant.
\end{abstract}

% ============================================================================
\section{Introduction}
\label{sec:intro}

The optimization viewpoint of machine learning casts training as the
minimization of an empirical risk
$f(\theta) = \frac{1}{N}\sum_{i=1}^{N} \ell(\theta; x_i, y_i)$
over the parameter vector $\theta \in \R^d$
\citep{robbins1951stochastic,goodfellow2016deep}. For modern
architectures $d$ is in the millions, $f$ is non-convex, and the
gradient $\nabla f(\theta)$ provides only local information. A long
line of work, both theoretical and empirical, has established that the
critical-point structure of $f$ in high dimensions is dominated by
\emph{saddle points} rather than by local minima
\citep{dauphin2014saddle,choromanska2015loss}. At a strict saddle the
gradient vanishes, the Hessian has both positive and negative
eigenvalues, and a vanilla gradient method --- which uses only
$\nabla f$ --- can stall for many iterations while the loss remains
nearly constant.

Two families of solutions have been proposed.
\textbf{Adaptive methods} (Adam \citep{kingma2015adam}, RMSProp,
AdaGrad, etc.) rescale the gradient by per-parameter second-moment
estimates, which empirically allows escape from many ill-conditioned
regions but does not exploit the loss value at points other than the
iterate.
\textbf{Perturbation methods} inject randomness near critical points to
cross the negative-curvature directions of the saddle.
\citet{ge2015escaping} and \citet{jin2017escape} establish that
gradient descent with isotropic Gaussian noise added on stagnation
detects strict saddles in a polynomial number of steps and converges
to a second-order stationary point.
\citet{lee2016gradient} prove that gradient descent almost surely does
not converge to a strict saddle, but the time spent traversing the
saddle's neighborhood remains the practical bottleneck.

\citet{vahedi2024spgd} introduce \emph{Steepest Perturbed Gradient
Descent} (SPGD), which differs from PGD-style noise injection in two
ways. First, SPGD perturbs uniformly (not via Gaussian noise tied to
the gradient norm). Second, and more importantly, SPGD draws \np{}
candidate perturbations and \emph{selects} the one yielding the
greatest loss decrease, fall\-ing back to the gradient step if no
candidate improves on the iterate. This combines a directed
gradient direction with a Monte Carlo exploration step. The original
work evaluates SPGD on Rastrigin, Ackley, Rosenbrock, and a 3-D
configuration-space engineering benchmark. The behavior of SPGD on
\emph{machine-learning} loss landscapes --- where the relevant geometry
is shaped by ReLU activations, batch normalization
\citep{ioffe2015batchnorm}, and stochastic mini-batching --- is to our
knowledge unstudied.

\paragraph{Contributions.}
\begin{enumerate}[leftmargin=*,itemsep=2pt,topsep=2pt]
  \item We provide an open implementation of SPGD as a PyTorch
        \verb|Optimizer| subclass alongside reference implementations
        of PGD and a same-mechanism random-perturbation control (RPGD).
  \item We design a controlled head-to-head comparison across four
        discriminating testbeds spanning synthetic, tabular, vision,
        and matrix-completion settings, each with a paired-seed
        protocol so optimizer-vs-optimizer effects are isolated from
        initialization noise; two further diagnostic experiments
        (Two Moons visualisation and the $\np \times \iterp$
        ablation) are reported in the appendix.
  \item We provide direct measurements of saddle-related behavior via
        stagnation episodes, in addition to the usual loss/accuracy
        endpoints, and report results that question the importance of
        ``saddle'' geometry in our deepest experiment.
  \item We isolate which mechanism drives SPGD's advantage by
        comparing SPGD against the random-perturbation control (RPGD).
        On CIFAR-10/ResNet-18 the gap is mean $+3.78$ pp test accuracy,
        paired across all seeds, supporting the claim that the
        steepest-selection rule (not perturbation alone) is the
        mechanism doing useful work.
  \item We characterize SPGD's sensitivity to its two
        hyper\-parameters \np{} and \iterp{} via a $4 \times 4$
        ablation, identifying a practical operating regime
        (Appendix~\ref{app:ablation}).
  \item We empirically test the SPGD paper's stated future-work
        hypothesis --- that the algorithm has ``potential to
        significantly enhance neural network training''
        \citep[Conclusion]{vahedi2024spgd} --- on a non-convex ML
        benchmark with documented saddle structure (MovieLens-100K
        matrix completion via Burer--Monteiro factorisation;
        \citealp{burer2003nonlinear,ge2016matrix,sun2018geometric}).
        Under fair experimental protocol, the paper-faithful SPGD
        selection rule fires at $30\%$ acceptance versus $9\%$ for
        the random-perturbation control --- direct evidence that the
        mechanism is empirically active --- but the held-out test
        RMSE gain over plain SGD is within seed-noise, splitting the
        paper's hypothesis into a mechanism-positive,
        generalisation-neutral finding.
\end{enumerate}

% ============================================================================
\section{Background}
\label{sec:background}

In high-dimensional non-convex losses the ratio of saddle points to
local minima grows exponentially in $d$, so the practical obstacle to
first-order optimisation is the slow traversal of saddle
neighbourhoods rather than convergence to a bad local minimum
\citep{dauphin2014saddle,choromanska2015loss}.  Perturbed gradient
descent (PGD) \citep{ge2015escaping,jin2017escape} addresses this by
adding Gaussian noise when $\|\nabla f\|$ is small; PGD finds a
second-order stationary point in $\tilde O(d/\varepsilon^{2})$
iterations w.h.p.

SPGD \citep{vahedi2024spgd} differs from PGD on three points:
(i) it perturbs unconditionally every \iterp{} steps, not on
stagnation detection; (ii) it draws \np{} candidates per perturbation
phase rather than one; and (iii) it commits the candidate with the
steepest loss decrease.  Algorithm~\ref{alg:spgd} restates the
procedure.

\begin{algorithm}[!htbp]
\caption{Steepest Perturbed Gradient Descent (SPGD)}
\label{alg:spgd}
\begin{algorithmic}[1]
\Require initial parameters $\theta_0$, learning rate $\eta$, perturbation amplitude $\amp$,
         number of candidates \np{}, perturbation interval \iterp{}, total steps $T$
\For{$t = 0, \ldots, T-1$}
  \If{$t \bmod \iterp{} = 0$ \textbf{and} $t > 0$}
    \State Sample \np{} candidates $\delta^{(k)} \sim \mathrm{Uniform}(-\amp, \amp)^d$
    \State Evaluate $f(\theta_t + \delta^{(k)})$ for $k = 1, \ldots, \np{}$
    \State Let $k^\star = \arg\min_{k} f(\theta_t + \delta^{(k)})$
    \If{$f(\theta_t + \delta^{(k^\star)}) \leq f(\theta_t)$}
      \State $\theta_t \gets \theta_t + \delta^{(k^\star)}$ \Comment{accept candidate (paper's $\leq$ rule)}
    \EndIf
  \EndIf
  \State $\theta_{t+1} \gets \theta_t - \eta \nabla f(\theta_t)$ \Comment{standard gradient step}
\EndFor
\end{algorithmic}
\end{algorithm}

% ============================================================================
\section{Methodology}
\label{sec:methodology}

\subsection{Optimizers under comparison}
\label{sec:optimizers}
We compare five algorithms, chosen so that each ablation isolates a
specific mechanism.

\begin{itemize}[leftmargin=*,itemsep=2pt,topsep=2pt]
  \item \textbf{SGD} \citep{robbins1951stochastic}: vanilla gradient
        descent, no momentum. The first-order baseline.
  \item \textbf{Adam} \citep{kingma2015adam}: adaptive per-parameter
        learning rate. The modern ML default; included to ground our
        results against current practice.
  \item \textbf{PGD} \citep{jin2017escape}: gradient descent with
        Gaussian noise added every \iterp{} steps; one perturbation,
        no selection.
  \item \textbf{RPGD} (\emph{random-perturbation control}): every
        \iterp{} steps draw \np{} candidates uniformly and accept
        \emph{one at random}. Same compute as SPGD; differs only in
        the selection rule.
  \item \textbf{SPGD} \citep{vahedi2024spgd}: as in
        Algorithm~\ref{alg:spgd}; the steepest of the \np{} candidates
        is selected.
\end{itemize}

\noindent The \textbf{SPGD-vs-RPGD} contrast is the cleanest test of
SPGD's central claim: both algorithms perturb the same way and pay the
same compute cost, so any difference in outcome can be attributed to
the selection rule. The \textbf{SPGD-vs-PGD} contrast tests whether
multi-candidate perturbation is helpful at all, and the
\textbf{SPGD-vs-SGD} contrast tests whether perturbation+selection
beats no perturbation.

\subsection{Diagnostics}
\label{sec:diagnostics}
Beyond the standard final-loss / final-accuracy endpoints we record
saddle-related behavior directly. A
\emph{stagnation episode} is a maximal run of consecutive steps with
$\|\nabla f(\theta_t)\|_2 < \varepsilon$. We log the number of
episodes, the longest episode, and the total number of stagnant steps.
We use the proposal's threshold $\varepsilon = 10^{-4}$ on the
synthetic benchmarks (Exp~\ref{sec:exp1}) and the Two Moons
diagnostic (App.~\ref{app:two-moons}) including its ablation
(App.~\ref{app:ablation}); $\varepsilon = 10^{-3}$ on OpenML
(Exp~\ref{sec:exp3}) and the CIFAR-10 mini\-batch gradient
(Exp~\ref{sec:exp4}), where the looser threshold accommodates
mini\-batch gradient noise. The proposal also listed an
\emph{escape-time} metric in addition to stagnation; we report
escape-time on Experiment~\ref{sec:exp6} (matrix completion), the
only setting where the iterate sits inside a documented saddle
plateau long enough for ``escape'' to have an objective definition.
Where stagnation episodes are zero across all runs of an experiment
(App.~\ref{app:two-moons}, Exp~\ref{sec:exp4},
App.~\ref{app:ablation}), escape-time is undefined and we say so
explicitly rather than reporting degenerate numbers.

\subsection{Paired-seed protocol}
\label{sec:paired}
For each (experiment, seed) the data split, parameter initialization,
and any per-seed randomness are fixed before the optimizer is
selected. The same five optimizers are then run on the
\emph{same} initial state. This converts an across-seed comparison
into a paired comparison: each seed produces a tuple of five outcomes
that share initialization. Paired differences are reported in the
results section because they are markedly less noisy than unpaired
means at our small seed counts ($S \in \{3, 5, 30\}$).

\subsection{Implementation}
The full software stack is PyTorch \citep{paszke2019pytorch} with
scikit-learn \citep{pedregosa2011sklearn} for synthetic-data
generation. Datasets from OpenML-CC18
\citep{vanschoren2014openml,bischl2021openmlcc18} are accessed via the
\texttt{openml} Python client. The CIFAR-10 anchor experiment uses
ResNet-18 \citep{he2016resnet} from \texttt{torchvision} with batch
normalization \citep{ioffe2015batchnorm} disabled in stages
\texttt{layer1} and \texttt{layer2} to deliberately preserve a
fraction of the saddle structure that BN would otherwise smooth out.
This concretises the proposal's ``no BN in early layers'' plan:
\texttt{layer1}/\texttt{layer2} are the ``early'' stages, and we
retained BN in \texttt{layer3}/\texttt{layer4} for training stability
at \texttt{lr=1e-2} on the gradient-only optimizers.
Source code, configurations, exact hyper\-parameters, and analysis
scripts are organized under \texttt{src/spgd\_study/} and
\texttt{experiments/}, and are publicly available at
\url{https://github.com/shyampatadia/spgd_research}.

% ============================================================================
\section{Experiments and Results}
\label{sec:results}

We present four discriminating experiments, increasing in scale from
$d=2$ analytic benchmarks to a $1.1\times 10^{7}$-parameter image
classifier and a $1.3\times 10^{4}$-parameter matrix-completion
problem; two further diagnostic experiments are reported in the
appendix.

% --------------------------------------------------------------------------
\subsection{Experiment 1: Non-Convex Benchmark Functions}
\label{sec:exp1}

\paragraph{Setup.} We run all five optimizers on Rastrigin, Ackley,
and Rosenbrock at $d \in \{2, 10, 50\}$, with 30 random initializations
per (function, dimension, optimizer). Convergence is declared at
$f(\theta) < \varepsilon = 10^{-2}$ from the global optimum. We
measure (i) the final function value, (ii) the iteration of first
convergence, and (iii) the success rate (fraction of seeds reaching
the threshold).

\paragraph{Results.} Figure~\ref{fig:exp1-trajectories} shows 2D
trajectories on each benchmark; SPGD and PGD escape multiple Rastrigin
basins that trap vanilla SGD. Figure~\ref{fig:exp1-convergence} plots
running-min loss vs.\ iteration for $d \in \{2,10,50\}$, and
Figure~\ref{fig:exp1-success} reports per-(function, dimension)
success rates. The cleanest signal appears on Ackley at $d=10$, where
SPGD's mean final loss is below RPGD's (paired delta of $-2.57$ in
mean loss, $-2.70$ in median). At $d=50$ all perturbation methods
struggle: a constant amplitude $\amp{} = 0.5$ that worked at $d=2$
becomes a $\sqrt{50}/\sqrt{2} \approx 5\times$ larger total
displacement and starts to produce destructive perturbations.
Per-parameter scaling of $\amp$, used in our subsequent experiments,
addresses this.

\begin{figure}[!htbp]
  \centering
  \includegraphics[width=\linewidth]{figures/exp1_trajectories_2d.png}
  \caption{2D trajectories of all five optimizers on Rastrigin, Ackley,
           and Rosenbrock. Stars mark the global optimum; thirty
           seeds per cell. SPGD's perturbation phases are visible as
           short ``hops'' that allow escape from local minima which
           vanilla gradient methods cannot leave.}
  \label{fig:exp1-trajectories}
\end{figure}

\begin{figure}[!htbp]
  \centering
  \includegraphics[width=\linewidth]{figures/exp1_convergence.png}
  \caption{Running-minimum loss vs.\ iteration on Experiment 1 (mean
           $\pm 1$ std across 30 seeds). Rows: dimension; columns:
           function. Symlog $y$-axis. The benefit of perturbation
           methods grows with dimension on Ackley but vanishes on
           Rosenbrock, which has no local minima.}
  \label{fig:exp1-convergence}
\end{figure}

\begin{figure}[!htbp]
  \centering
  \includegraphics[width=0.95\linewidth]{figures/exp1_success_rate.png}
  \caption{Success rate (fraction of seeds reaching within
           $\varepsilon = 10^{-2}$ of the global optimum) per
           (function, dimension, optimizer). SPGD matches or exceeds
           RPGD on every Rastrigin and Ackley cell.}
  \label{fig:exp1-success}
\end{figure}
\FloatBarrier

% --------------------------------------------------------------------------
% Experiment 2 (Two Moons) is a visualization-only diagnostic and has
% been moved to Appendix~\ref{app:two-moons}.

\subsection{Experiment 3: OpenML-CC18 Tabular Datasets}
\label{sec:exp3}

\paragraph{Setup.} We select five OpenML-CC18 classification datasets
spanning a range of sizes and feature types: \texttt{credit-g}
($n{=}1000$, $p{=}20$), \texttt{kc1}, \texttt{spambase},
\texttt{segment}, \texttt{vehicle}. A 3-layer MLP is trained on each
for all five optimizers across three seeds (75 runs total). All
datasets are normalized; categorical features are one-hot encoded.

\paragraph{Results.} Figure~\ref{fig:exp3-acc} shows test accuracy per
(dataset, optimizer) and Figure~\ref{fig:exp3-loss} shows training-loss
trajectories in two complementary views: the left panel includes all
five optimizers on log-$y$ to make Adam's dominance visible, and the
right panel shows the four gradient-only methods on linear-$y$ to
resolve the SGD/PGD/RPGD/SPGD differences that the log-$y$ view
compresses into a single thick band. The two panels show the same
data; we plot both so the gradient-only zoom is not read as Adam being
hidden. The pattern from Experiment~\ref{sec:exp1}
reproduces: SPGD's mean test accuracy is above RPGD's on a majority of
(dataset, seed) pairs, supporting the hypothesis that the
steepest-selection rule transfers from synthetic to real heterogeneous
tabular data. Adam dominates absolute accuracy on all five datasets,
consistent with the observation that on small tabular tasks a
well-conditioned adaptive method has little incentive to perturb.
Stagnation episodes ($\|\nabla f\|_2 < 10^{-3}$) were rare on this
suite: only \texttt{credit-g} and \texttt{vehicle} produced any
non-zero counts, and even there the totals were a handful of episodes
per run, so we omit the per-dataset stagnation figure as
uninformative and report the diagnostic in text only.

\begin{figure}[!htbp]
  \centering
  \includegraphics[width=\linewidth]{figures/exp3_acc_per_dataset.png}
  \caption{OpenML-CC18 test accuracy per dataset (mean $\pm 1$ std
           across 3 seeds). Adam consistently leads; SPGD is at or
           above SGD on 4/5 datasets and above RPGD on a majority of
           paired comparisons.}
  \label{fig:exp3-acc}
\end{figure}

\begin{figure}[!htbp]
  \centering
  \begin{subfigure}{\linewidth}
    \centering
    \includegraphics[width=\linewidth]{figures/exp3_loss_curves.png}
    \caption{All five optimisers, log-$y$. Adam's curve is the lowest
             on all five datasets.}
    \label{fig:exp3-loss-all}
  \end{subfigure}

  \vspace{0.6em}

  \begin{subfigure}{\linewidth}
    \centering
    \includegraphics[width=\linewidth]{figures/exp3_loss_curves_grad_only.png}
    \caption{Gradient-only methods, linear-$y$ (Adam excluded for
             readability). The linear axis exposes SPGD's modest but
             consistent edge over RPGD on plateaus SGD also fails
             to leave.}
    \label{fig:exp3-loss-grad}
  \end{subfigure}
  \caption{Training-loss curves on OpenML-CC18 (mean $\pm 1$ std,
           3~seeds). Both panels plot the same runs; the bottom panel
           is a y-axis zoom on the gradient-only band that the top
           panel compresses.}
  \label{fig:exp3-loss}
\end{figure}

\FloatBarrier

% --------------------------------------------------------------------------
\subsection{Experiment 4: CIFAR-10 / ResNet-18 (Anchor Experiment)}
\label{sec:exp4}

\paragraph{Setup.} ResNet-18 \citep{he2016resnet} with batch
normalization disabled in \texttt{layer1} and \texttt{layer2} (to
preserve saddle structure that fully-BN networks tend to smooth out)
is trained on CIFAR-10 \citep{krizhevsky2009cifar} for 50 epochs
across all five optimizers and three seeds (15 runs).
Hyperparameters: \texttt{lr=1e-2} for the gradient-only optimizers,
\texttt{lr=1e-3} for Adam, $\amp = 10^{-3}$ (chosen as $\sim\!2\%$ of
typical Kaiming-initialized weight magnitude
\citep{he2015delving}), $\np = 5$, $\iterp = 200$ minibatch steps
($\sim\!2$ perturbation phases per epoch). Runs executed as a Slurm
array on two NVIDIA A30 GPUs (mean wall time $\sim 7.6$ min/run).

\paragraph{Results.} Table~\ref{tab:exp4-summary} reports aggregate
metrics; Figure~\ref{fig:exp4-loss} shows training-loss trajectories
on log-$y$, and Figure~\ref{fig:exp4-paired} reports the
\emph{paired} per-seed deltas SPGD$-$RPGD and SPGD$-$SGD. The
headline numbers:

\begin{itemize}[leftmargin=*,itemsep=1pt,topsep=2pt]
  \item Mean test accuracies, three seeds, 50 epochs:
        Adam $0.9015 \pm 0.014$,
        SPGD $0.8277 \pm 0.028$,
        SGD  $0.8268 \pm 0.019$,
        PGD  $0.8091 \pm 0.014$,
        RPGD $0.7899 \pm 0.046$.
  \item \textbf{SPGD beats RPGD on all three seeds} with paired mean
        $+3.78$ pp test accuracy (per-seed: $+4.42$, $+1.14$, $+5.77$
        pp). This is the single clearest piece of evidence in the
        study that SPGD's steepest-selection rule contributes signal
        beyond random perturbation.
  \item SPGD ties SGD overall (paired mean $+0.09$ pp), consistent
        with the next finding.
  \item \textbf{Stagnation episodes were zero for all 15 runs}. None
        of the five optimizers spent any consecutive step with
        $\|\nabla f\| < 10^{-3}$. The proposal's intervention
        (disabling BN in \texttt{layer1}+\texttt{layer2}) was
        \emph{insufficient} to expose saddle/plateau structure on
        this CIFAR-10/ResNet-18 setup --- the surviving normalization
        in \texttt{layer3} and \texttt{layer4} keeps minibatch
        gradient norms above the threshold throughout training.
  \item Adam's adaptive step is the strongest absolute optimizer
        ($\sim\!90\%$); the proposal's acceptance criterion of
        ``SGD $\geq 0.90$'' was \emph{not} met (SGD: 0.83), an
        expected consequence of the deliberate handicaps used to
        preserve saddle structure.
\end{itemize}

\begin{table}[!htbp]
\centering
\small
\begin{tabular}{lcccr}
\toprule
\textbf{Optimizer} & \textbf{Test acc.\ (mean $\pm$ std)} & \textbf{Final train loss} & \textbf{Stagnation eps.} & \textbf{Wall (s)} \\
\midrule
SGD   & $0.8268 \pm 0.019$ & $0.249 \pm 0.044$ & $0$ & $451$ \\
Adam  & $0.9015 \pm 0.014$ & $0.119 \pm 0.079$ & $0$ & $464$ \\
PGD   & $0.8091 \pm 0.014$ & $0.311 \pm 0.015$ & $0$ & $453$ \\
RPGD  & $0.7899 \pm 0.046$ & $0.337 \pm 0.079$ & $0$ & $458$ \\
SPGD  & $0.8277 \pm 0.028$ & $0.293 \pm 0.046$ & $0$ & $460$ \\
\bottomrule
\end{tabular}
\caption{Experiment 4 aggregate results. CIFAR-10, ResNet-18 with BN
         disabled in \texttt{layer1}+\texttt{layer2}, 50 epochs,
         3 seeds. All differences are with respect to the same
         per-seed initialization (paired protocol).}
\label{tab:exp4-summary}
\end{table}
\FloatBarrier

\begin{figure}[!htbp]
  \centering
  \begin{subfigure}[t]{0.49\linewidth}
    \includegraphics[width=\linewidth]{figures/exp4_loss_curves.png}
    \caption{All five optimizers (log $y$).}
    \label{fig:exp4-loss-all}
  \end{subfigure}\hfill
  \begin{subfigure}[t]{0.49\linewidth}
    \includegraphics[width=\linewidth]{figures/exp4_loss_curves_grad_only.png}
    \caption{Gradient-only methods (linear $y$, Adam excluded).}
    \label{fig:exp4-loss-grad}
  \end{subfigure}
  \caption{CIFAR-10/ResNet-18 training cross-entropy (mean $\pm 1$ std
           across 3 seeds). Both panels plot the same runs. The log-$y$
           view (left) includes all five optimizers and shows Adam's
           dominance from epoch~1; the linear-$y$ view (right) is the
           same plot with Adam removed and the axis rescaled, so the
           SGD/PGD/RPGD/SPGD differences --- which the log-$y$ view
           compresses into a single thick band near the axis --- become
           visible. The right panel is a presentation choice for
           readability, not a data exclusion.}
  \label{fig:exp4-loss}
\end{figure}

\begin{figure}[!htbp]
  \centering
  \includegraphics[width=0.85\linewidth]{figures/exp4_paired_spgd.png}
  \caption{\textbf{Key paired comparison.} Per-seed test-accuracy
           deltas SPGD$-$RPGD (red) and SPGD$-$SGD (blue) on CIFAR-10.
           Dashed lines mark the per-comparison means. SPGD beats
           RPGD on every seed (mean $+3.78$ pp), the cleanest
           evidence in the study that the \emph{selection rule} ---
           not perturbation alone --- drives the improvement.}
  \label{fig:exp4-paired}
\end{figure}
\FloatBarrier

% --------------------------------------------------------------------------
% Experiment 5 ($N_P \times \mathit{IterP}$ ablation on Two Moons) is a
% non-discriminating diagnostic on a benign landscape and has been moved
% to Appendix~\ref{app:ablation}.

\subsection{Experiment 6: Matrix Completion on MovieLens-100K}
\label{sec:exp6}

The preceding experiments (and the appendix diagnostics) evaluate
SPGD as a neural-network trainer.  This section directly tests the SPGD paper's
stated future-work hypothesis --- that SPGD has ``potential to
significantly enhance neural network training'' through ``improving
convergence rates and navigating complex parameter spaces''
\citep[Conclusion]{vahedi2024spgd} --- on a non-convex ML benchmark
with documented saddle structure: low-rank Burer--Monteiro
factorisation \citep{burer2003nonlinear} of the MovieLens-100K rating
matrix \citep{harper2015movielens}.  The optimisation literature
establishes that small-init Burer--Monteiro factorisation has a saddle
near the origin that gradient descent escapes only slowly
\citep{ge2016matrix,sun2018geometric}, exactly the regime SPGD's
selection rule is designed for.

\paragraph{Setup.}  We complete the rating matrix
$M \in \mathbb{R}^{943 \times 1682}$ via $\widehat{M} = U V^{\top}$
with $U \in \mathbb{R}^{943 \times 5}$ and
$V \in \mathbb{R}^{1682 \times 5}$, minimising the L2-regularised
mean-squared reconstruction error
\begin{equation*}
  \mathcal{L}(U, V) \;=\;
    \frac{1}{|\Omega|} \sum_{(i,j) \in \Omega}
      \bigl( (UV^{\top})_{ij} - r_{ij} \bigr)^{2}
    \;+\; \lambda
      \bigl( \overline{\|U\|^2} + \overline{\|V\|^2} \bigr)
\end{equation*}
on a random 80/20 train/test split of the $100{,}000$ observed
entries, with $\lambda = 0.05$ applied identically to all four
methods (a fairness control: without it, Adam overfits while
SGD/SPGD enjoy implicit early-stopping regularisation).  The factors
are initialised i.i.d.\ from $\mathcal{N}(0, 0.01^2)$, placing the
iterate inside the saddle-rich neighbourhood of the origin.  We
compare SGD, Adam, RPGD, and SPGD over $T{=}1500$ full-batch gradient
evaluations across three seeds.  SPGD and RPGD use the verified
optimiser classes from \texttt{src/spgd\_study/optimizers/},
faithful to Algorithm~\ref{alg:spgd}: $\amp = 0.01$ per coordinate,
Uniform$(-\amp, +\amp)$ candidates, $\np{=}8$, $\iterp{=}5$, and the
paper's $\le$ acceptance rule (a candidate is committed only if its
loss is no greater than that of the current iterate).

\paragraph{Aggregate results.}  Table~\ref{tab:exp6-summary} reports
the aggregate metrics.  All four methods land within $0.02$ test
RMSE of one another.  Paired per-seed deltas
(Figure~\ref{fig:exp6-paired}) place SPGD slightly ahead of SGD on
training MSE (paired $\Delta = -0.0034$) but indistinguishable from
RPGD on test RMSE ($\Delta = +0.0003$, sign mixed across seeds).

\paragraph{Escape-time diagnostic (proposal Objective~3).}
We define escape-time as the first step at which the training MSE has
fallen $5\%$ below its initial value, the only setting in this study
where the iterate sits inside a documented saddle plateau long enough
for ``escape'' to have an objective definition (the proposal's
escape-time is undefined in App.~\ref{app:two-moons},
Exp~\ref{sec:exp4}, and App.~\ref{app:ablation}, where stagnation
episodes are zero).  Per-method
mean escape steps are
\textbf{Adam $40 \pm 17$, SPGD $320 \pm 17$, RPGD $360 \pm 30$, SGD
$370 \pm 17$}; paired per-seed deltas are
\textbf{$\mathrm{SPGD} - \mathrm{SGD} = -50$ steps mean (per-seed:
$-60, -30, -60$)} and
\textbf{$\mathrm{SPGD} - \mathrm{RPGD} = -40$ steps mean (per-seed:
$-30, -30, -60$)}: SPGD escapes the saddle plateau on \emph{every
seed} ahead of both same-data baselines, but the lead is consumed by
the time the methods converge.  Numbers are reproducible from
\texttt{results/exp6\_escape\_time\_paired.json}.

\paragraph{The acceptance-rate diagnostic.}
The most informative number in this experiment is SPGD's
perturbation acceptance rate:
\textbf{$30 \pm 2\%$ across three seeds (88--97 of 299 perturbation
phases produced a strictly improving candidate)}.  The matched random
control RPGD accepts only $\bm{9 \pm 1\%}$ of phases, since a
uniformly random pick from $\np{=}8$ candidates rarely lands on the
argmin and rarely satisfies the $\le$ rule.  The threefold gap is
direct evidence that SPGD's selection rule is firing as designed on a
$13{,}125$-parameter non-convex ML problem --- the failure mode of a
``silent'' SPGD that never accepts is ruled out by this diagnostic.

\begin{table}[!htbp]
\centering
\small
\begin{tabular}{lcccc}
\toprule
\textbf{Method} & \textbf{Final train MSE} & \textbf{Test RMSE} &
\textbf{Acceptance} & \textbf{Wall (s)} \\
\midrule
SGD   & $0.650 \pm 0.007$ & $0.949 \pm 0.007$ & ---           & $2.5$ \\
Adam  & $0.596 \pm 0.002$ & $0.963 \pm 0.003$ & ---           & $2.8$ \\
RPGD  & $0.649 \pm 0.002$ & $0.947 \pm 0.002$ & $9 \pm 1\%$   & $4.2$ \\
SPGD  & $0.647 \pm 0.005$ & $0.948 \pm 0.002$ & $30 \pm 2\%$  & $4.2$ \\
\bottomrule
\end{tabular}
\caption{Experiment~6 aggregate results.  MovieLens-100K rank-5
         Burer--Monteiro matrix completion, $\lambda{=}0.05$, $1500$
         full-batch steps, three seeds.  Acceptance is the fraction
         of SPGD/RPGD perturbation phases that committed a candidate
         (i.e.\ found a candidate satisfying the $\le$ acceptance
         rule).  All four methods see the identical L2-regularised
         loss inside their respective forward passes.}
\label{tab:exp6-summary}
\end{table}

\begin{figure}[!htbp]
\centering
\includegraphics[width=\linewidth]{figures/exp6_mc_curves.png}
\caption{Experiment~6 --- training MSE (left, log scale) and held-out
         test RMSE (right) over $1500$ full-batch steps on
         MovieLens-100K with rank-5 Burer--Monteiro factors initialised
         at scale $0.01$ (mean $\pm 1$ std across three seeds).
         SPGD's training loss reaches the $5\%$-drop escape threshold
         a paired mean of $50$ steps before SGD and $40$ steps before
         RPGD (3/3 seeds; see ``Escape-time diagnostic'' above); the
         lead largely dissipates by the end of training.  Adam's training loss is lower throughout
         but its test RMSE plateaus higher --- a classic mild-overfit
         signature, despite identical L2 regularisation.}
\label{fig:exp6-curves}
\end{figure}

\begin{figure}[!htbp]
\centering
\includegraphics[width=0.85\linewidth]{figures/exp6_mc_paired.png}
\caption{Experiment~6 --- per-seed paired deltas of held-out test
         RMSE between SPGD and the three baselines.  Negative bars
         indicate SPGD ahead of the baseline on that seed.  SPGD beats
         SGD on $2/3$ seeds (mean $\Delta = -0.0018$), is
         statistically tied with RPGD ($\Delta = +0.0003$, sign mixed
         across seeds), and beats Adam on $3/3$ seeds
         ($\Delta = -0.0159$) --- the last margin is dominated by
         Adam's mild overfit, not by SPGD's optimisation advantage.}
\label{fig:exp6-paired}
\end{figure}

\paragraph{Interpretation.}
SPGD's mechanism is empirically active on this benchmark --- the
$30\%$ vs.\ $9\%$ acceptance gap is unambiguous evidence that the
selection rule fires more often than chance and that the algorithm is
running as the paper specifies.  However, the optimisation advantage
this confers (a $50$-step paired-mean head-start over SGD, $40$ over
RPGD, on the $5\%$-drop escape-time metric, 3/3 seeds) does
\emph{not} translate into a measurable held-out generalisation gain: the test-RMSE gap to plain SGD is
$-0.0018 \pm 0.005$ with the per-seed sign mixed, and the gap to RPGD
is essentially zero.  Reading the paper's future-work claim
empirically, our answer is split: the selection rule's mechanism
\emph{does} generalise from 2D test functions to a high-dimensional
non-convex ML problem (mechanism-active), but the ``significant
enhancement'' the paper hypothesises does not materialise on this
benchmark (effect-small).
\FloatBarrier

% ============================================================================
\section{Conclusion}
\label{sec:conclusion}

This work is the first systematic empirical study of SPGD on
machine-learning loss landscapes.  Across four discriminating
testbeds --- analytic benchmarks ($d \in \{2,10,50\}$), five
OpenML-CC18 tabular datasets, a CIFAR-10/ResNet-18 anchor experiment,
and MovieLens-100K matrix completion --- SPGD's steepest-selection
rule provides a measurable advantage over its same-mechanism
random-perturbation control (RPGD).  The advantage is cleanest on
CIFAR-10/ResNet-18: SPGD beats RPGD on all three seeds with a
paired-mean delta of $+3.78$ percentage points test accuracy.  On
matrix completion --- which directly tests the SPGD paper's stated
NN-training hypothesis \citep[Conclusion]{vahedi2024spgd} --- the
selection rule fires at threefold the rate of the random control
($30\%$ vs.\ $9\%$ acceptance) and SPGD escapes the saddle plateau
$50$ steps before SGD on all three seeds, but the held-out RMSE
gap is within seed noise: mechanism-positive, generalisation-neutral.

\paragraph{Caveats.}  No optimiser in our study exhibited sustained
gradient stagnation at $\varepsilon = 10^{-3}$, so SPGD's CIFAR-10
advantage cannot be attributed to plateau escape per se and likely
reflects a more subtle benefit of biased candidate selection.
Adam dominates absolute accuracy throughout because its adaptive
preconditioning is structurally different; the correct head-to-head
comparison is SPGD vs.\ same-compute random control, not SPGD vs.\
Adam.  Three seeds at the CIFAR-10 scale is small, mitigated only
by the paired protocol.

\paragraph{Recommendation.}  Future perturbation-based optimisers
should be benchmarked against a random-selection control of equal
compute --- a comparison missing from much of the existing literature
--- since on our deepest setting it is the \emph{selection rule},
not perturbation alone, that does the useful work.

% ============================================================================
\clearpage
\appendix
\renewcommand{\thesection}{\Alph{section}}

\begin{center}
{\Large\bfseries Appendix}
\end{center}

\noindent The two appendix experiments are diagnostic only: Two Moons
is a visualisation sanity check (the landscape is too benign to
discriminate optimisers), and the $\np \times \iterp$ ablation runs
on the same Two Moons setup so its accuracy heatmap is also
non-discriminating.  We retain them for transparency with the
proposal and because the figures expose secondary phenomena
(geometric difference between Adam and the gradient methods;
SPGD's compute--accuracy trade-off curve) that would otherwise be
unsupported assertions in the main text.

\vspace{0.6em}

\section{Two Moons --- Visualisation Sanity Check (proposal Exp~2)}
\label{app:two-moons}
\label{sec:exp2}

This experiment is a \emph{visualisation-only} diagnostic.  The
dataset is benign enough that a $2 \to 32 \to 32 \to 1$ MLP solves it
from any reasonable initialisation, so the saddle-escape mechanism
SPGD targets does not engage and the optimisers' final test
accuracies cluster within seed noise.  We retain the experiment in the
appendix because (i) it verifies the paired-seed protocol on a small
example, (ii) it reproduces \citet{vahedi2024spgd}'s
trajectory-visualisation style on a neural network rather than a 2-D
test function, and (iii) the resulting picture exposes a real
geometric difference between adaptive and non-adaptive optimisers in
weight space.  The head-to-head numbers reported in the main paper
come from Exps~\ref{sec:exp1}, \ref{sec:exp4}, and~\ref{sec:exp6}.

\paragraph{Setup.} A $2 \to 32 \to 32 \to 1$ MLP is trained on the
\texttt{make\_moons} dataset (1{,}000 points, noise $0.2$) full-batch
for 2{,}000 steps. Five optimisers, five seeds, 25 runs total.
Following \citet{vahedi2024spgd}'s visualisation style, seed 0 of each
optimiser additionally records its weight trajectory at every 10th
step; we render the loss landscape in a \emph{shared} top-2 PCA basis
(union of all five trajectories' displacements from $\theta_0$) with
each optimiser's trajectory overlaid.  Centering on $\theta_0$ and
sharing the basis means PC1 and PC2 mean the same thing in every
panel and the paired-seed initialisation lands at $(0,0)$ everywhere.

\paragraph{Results --- numbers.}  Final test accuracies
(mean $\pm 1$ std across 5 seeds): Adam $0.977 \pm 0.010$,
PGD $0.970 \pm 0.018$, SPGD $0.970 \pm 0.015$,
RPGD $0.967 \pm 0.019$, SGD $0.966 \pm 0.019$.  The four gradient
methods are within seed noise of each other; Adam wins for the
well-known reason that adaptive scaling fits a benign basin
quickly.  This experiment does not adjudicate SPGD vs.\ baselines.

\paragraph{Results --- trajectories.}
Figure~\ref{fig:exp2-landscape} is the visualisation showpiece.  Two
features are robust across seeds: (i) the four gradient methods (SGD,
PGD, RPGD, SPGD) all move predominantly along the shared PC2
direction, while Adam moves substantially along PC1 --- a direct
picture of adaptive scaling exploring directions the gradient methods
do not preferentially use; (ii) SPGD's path is faintly serrated
relative to SGD's, the small lateral excursions are SPGD's accepted
perturbation candidates.  These are diagnostic observations about
\emph{geometry}, not claims about \emph{which optimiser wins}.

Figure~\ref{fig:exp2-loss} shows training-loss curves and
Figure~\ref{fig:exp2-gradnorm} the gradient-norm history with the
proposal's stagnation threshold $\varepsilon = 10^{-4}$ drawn for
reference.  On this dataset the loss landscape is sufficiently benign
that \emph{none} of the five optimisers exhibits sustained stagnation
($\|\nabla f\| < 10^{-4}$ for $\geq 5$ consecutive steps); the
``saddle escape'' mechanism is therefore not observable in the
endpoint statistics here, only as the small trajectory feature noted
above.

\begin{figure}[!htbp]
  \centering
  \includegraphics[width=\linewidth]{figures/exp2_loss_landscape.png}
  \caption{Two Moons training-BCE loss surface in a \emph{shared}
           PCA basis (seed 0). PC1 and PC2 are the top-2 principal
           directions of the union of weight displacements
           $\theta(t) - \theta_0$ across all five optimisers, so every
           panel projects onto the same 2-D plane through 1185-D weight
           space. The colour map and white iso-loss contours show the
           training BCE loss evaluated at $\theta_0 + a\,\mathbf{u} +
           b\,\mathbf{v}$, clipped to the loss range the trajectories
           actually visit (yellow = approaching the clip ceiling,
           purple = low-loss basin). Contour levels are log-spaced so
           the basin's structure is legible; an unclipped colour scale
           would be dominated by the saturation strip on the left edge,
           where weights at the boundary of the sweep grid push the
           untrained-network BCE loss to $\sim$50. Green circle =
           shared initialisation $\theta_0$ at the origin, red ``X'' =
           final point, white curve = weight trajectory.  The geometric
           reading: \emph{the four gradient methods all move
           predominantly along PC2}, whereas \emph{Adam's trajectory
           has substantial PC1 motion} --- per-parameter adaptive
           scaling moves the iterate along a direction the gradient
           methods do not preferentially explore. SPGD's white path is
           faintly serrated relative to SGD's: those tiny lateral
           excursions are SPGD's perturbation phases.}
  \label{fig:exp2-landscape}
\end{figure}

\begin{figure}[!htbp]
  \centering
  \begin{subfigure}[t]{0.49\linewidth}
    \includegraphics[width=\linewidth]{figures/exp2_loss_curves.png}
    \caption{Training BCE loss (mean $\pm 1$ std across 5 seeds, log scale).}
    \label{fig:exp2-loss}
  \end{subfigure}\hfill
  \begin{subfigure}[t]{0.49\linewidth}
    \includegraphics[width=\linewidth]{figures/exp2_grad_norm_history.png}
    \caption{Gradient-norm trajectories with the proposal's stagnation
             threshold $\varepsilon = 10^{-4}$ drawn as a dashed line.}
    \label{fig:exp2-gradnorm}
  \end{subfigure}
  \caption{Two Moons summary. Left: training loss; Adam pulls away
           early. Right: $\|\nabla f\|_2$ vs.\ training step on log-$y$
           --- no optimiser's gradient norm dips below the stagnation
           threshold, which is why the stagnation episode count is
           zero across all 25 runs.}
  \label{fig:exp2-summary}
\end{figure}
\FloatBarrier

\section{$\np \times \iterp$ Hyperparameter Ablation (proposal Exp~5)}
\label{app:ablation}
\label{sec:exp5}

This appendix reports SPGD's sensitivity to its two perturbation
hyperparameters.  The ablation runs on the same Two Moons setup as
Appendix~\ref{app:two-moons}.  Because that landscape is
non-discriminating (zero stagnation across all 25 main-experiment
runs) the ablation similarly does not produce a saddle-escape
signal --- every cell of the grid sees zero stagnation episodes ---
and is consequently a hyperparameter-robustness diagnostic rather
than a discriminating benchmark.  We retain it because the
\emph{compute--accuracy trade-off} it reveals is independent of
saddle structure.

\paragraph{Setup.} Grid: $\np \in \{1, 5, 10, 20\}$ and
$\iterp \in \{5, 10, 50, 100\}$, five seeds per cell, totalling 80
runs.  Total candidate forward passes per run
$= (n_{\text{steps}} / \iterp{}) \cdot \np{}$, which spans roughly
$20$ to $8000$ across the grid (a $\sim 400\times$ compute range
with all other hyper\-parameters fixed).

\paragraph{Results.} Figure~\ref{fig:exp5-acc-loss} reports the mean
test accuracy and final-training-loss heatmaps;
Figure~\ref{fig:exp5-walltime} shows the corresponding wall-time
heatmap.  Mean stagnation episodes are zero across the entire grid,
so the stagnation heatmap is omitted as uninformative.  The accuracy
heatmap is largely flat: SPGD on Two Moons is robust to the precise
choice of \np{} and \iterp{} within reasonable bounds.  The
final-loss heatmap shows a faint trend toward smaller \iterp{} (more
frequent perturbation phases) reaching slightly lower training loss
without a corresponding accuracy gain, consistent with SPGD finding
broader basins rather than deeper ones.  The cheapest cell within
$1\%$ of the best mean test accuracy identifies a practical operating
regime that trades essentially no accuracy for a substantial reduction
in candidate forward-pass cost.  The wall-time heatmap follows the
predicted $(n_{\text{steps}}/\iterp{}) \cdot \np{}$ scaling, so when
SPGD is deployed in a compute-budgeted setting, choosing $\np{}$ and
$\iterp{}$ at the cheap-but-not-worst-accuracy corner of the heatmap
is the natural recommendation.

\begin{figure}[!htbp]
  \centering
  \begin{subfigure}[t]{0.49\linewidth}
    \includegraphics[width=\linewidth]{figures/exp5_heatmap_acc.png}
    \caption{Mean test accuracy.}
    \label{fig:exp5-acc}
  \end{subfigure}\hfill
  \begin{subfigure}[t]{0.49\linewidth}
    \includegraphics[width=\linewidth]{figures/exp5_heatmap_final_loss.png}
    \caption{Mean final training loss.}
    \label{fig:exp5-loss}
  \end{subfigure}
  \caption{Ablation heatmaps. Left: test accuracy is largely
           insensitive to \np{} and \iterp{} within the explored
           range, suggesting SPGD's hyper\-parameter choice is
           forgiving on benign landscapes. Right: final BCE training
           loss exposes a small but visible compute-quality trade-off
           --- smaller \iterp{} reaches slightly lower training loss
           --- which the accuracy view alone cannot resolve.}
  \label{fig:exp5-acc-loss}
\end{figure}

\begin{figure}[!htbp]
  \centering
  \includegraphics[width=0.55\linewidth]{figures/exp5_heatmap_walltime.png}
  \caption{Per-run wall time across the $\np \times \iterp$ grid. The
           pattern follows the analytic prediction
           $\propto (n_{\text{steps}}/\iterp{}) \cdot \np{}$. Choosing
           the smallest \np{} and largest \iterp{} that does not
           sacrifice accuracy is the practical recipe for cost-
           effective SPGD.}
  \label{fig:exp5-walltime}
\end{figure}
\FloatBarrier

% ============================================================================
\section*{Acknowledgments}
The author thanks the WPI Academic Computing team for access to the
Turing cluster, on which the CIFAR-10 anchor experiment was carried
out. This work was completed in fulfillment of the final-project
requirement of MA 551 (Computational Statistics) at Worcester
Polytechnic Institute.

% ---------- Bibliography ----------------------------------------------------
\bibliographystyle{plainnat}
\bibliography{references}

\end{document}
